\begin{dissertationepilogue}

    % Fix a bug in the numbering of the epilogue chapters?
    \setcounter{chapter}{4}
    \setcounter{section}{0}

    \section{Conclusion}
    Insulin resistance (IR) is a global epidemic, and it is imperative to shift from broad-spectrum pharmacology toward precision molecular targeting to combat this multifaceted condition. While the therapeutic potential of PPAR$\gamma$ has been recognized for decades, the clinical utility of thiazolidinediones (TZDs) remains constrained by their adverse side-effect profile\cite{Yki-Jarvinen2004-bc}. As a master regulator of adipocyte differentiation, PPAR$\gamma$ orchestrates the development and maintenance of adipose tissue\cite{Tontonoz2008-pe} and contributes to distinctions between adipose depots: adipocytes derived from visceral adipose tissue, which is associated with exacerbating the pathogenesis of IR-related diseases\cite{Preis2010-eh,Agrawal2023-ee}, exhibit reduced PPARG activity compared to subcutaneous adipocytes\cite{Bradford2019-ps,Mathur2022-fx}. Despite these insights, the molecular mechanisms by which depot specific adipocytes affect adipose tissue dysfunction and systemic insulin resistance remain poorly defined\cite{Sebo2019-nh,Sakers2022-wb}. This dissertation addressed these challenges by integrating functional genomics with human genetics to uncover a novel translational start site within PPARG and by developing a computational infrastructure to generate and interrogate depot-specific adipocytes.

    
    %%%%%%%%%%%%%%%%%%%%%%%%%%%%%%%%%%%%%%%%%%%%%%%%%%%%%%%%%%%%%%%%%%%%%%%%%%%%
    \section{Summary}
    %%%%%%%%%%%%%%%%%%%%%%%%%%%%%%%%%%%%%%%%%%%%%%%%%%%%%%%%%%%%%%%%%%%%%%%%%%%%

    In Chapter 1, we investigate the coding landscape of PPARG and identify an alternative translational start site at methionine 135. This discovery redefines our understanding of the PPARG locus, demonstrating that the upstream region encoding the activation-function 1 (AF-1) domain of the protein is not necessary for function. The novel translational initiation site generates the isoform PPAR$\gamma$ M135, which manifests superior transactivation of target genes than the canonical PPAR$\gamma$2 isoform in response to TZD treatment.
    
    The therapeutic relevance of PPAR$\gamma$ M135 was confirmed in Chapter 2 through a combination of in vitro modeling and human genetic association studies. SGBS human adipocytes engineered to express PPAR$\gamma$ M135 exhibited increased insulin-sensitizing capacity relative to wild-type controls. Consistent with the heightened activity of this truncated isoform, human genetic variants that attenuate the PPARG AF-1 domain correlate with favorable metabolic profiles, as quantified by a metabolic syndrome severity score (METSS) developed within this study. Together, these findings suggest that AF-1 inhibition or selective expression of PPAR$\gamma$ M135 may amplify the therapeutic benefits of PPAR$\gamma$ agonism while minimizing the adverse effects associated with conventional TZDs.
    
    Finally, in Chapter 3, we provide a framework for modeling depot-specific adipocytes. SGBS cells were verified as a suitable progenitor model for both visceral and subcutaneous adipocytes, though other multipotent stem cells or fibroblasts may offer similar utility. We identified 46 candidate transcription factors (TFs) that distinguish between subcutaneous and visceral (pre)adipocytes and applied orthogonal methods to substantiate 17 TFs, including PPARG. To evaluate the reprogrammed cells, we proposed a computational blueprint that implemented cell identity scoring methods on pseudobulked single cell RNAseq and achieved high accuracy. This work provides a foundation to further study the mechanisms contributing to IR in the cellular contexts in which they manifest.

    %%%%%%%%%%%%%%%%%%%%%%%%%%%%%%%%%%%%%%%%%%%%%%%%%%%%%%%%%%%%%%%%%%%%%%%%%%%%
    \section{Limitations and Future Directions}
    %%%%%%%%%%%%%%%%%%%%%%%%%%%%%%%%%%%%%%%%%%%%%%%%%%%%%%%%%%%%%%%%%%%%%%%%%%%%

    The PPAR$\gamma$ M135 isoform demonstrates that extensively studied genes still harbor critical, undiscovered regulatory mechanisms with profound clinical implications. The plethora of data generated from single cell functional genomics methods could be helpful in this pursuit.
    
    One innovation in this thesis was the development of the metabolic syndrome severity score (METSS), which assessed metabolic syndrome in the UK Biobank based on five clinical components. However, this score had limited application and scope, as the clinical landscape progresses quickly, especially with multifaceted conditions. This was recently exemplified by the evolution of steatotic liver disease nomenclature. In 2021, I wrote a review about the human genetics underlying non-alcoholic fatty liver disease (NAFLD), in which I discussed the difficulties of diagnosing and classifying NAFLD and annotated each known genetic variant by the specific phenotypes it was found to be associated with, such as histologically defined steatosis, liver fat measured by imaging modalities, and liver enzyme levels\cite{Du2021-bq}. In the review process, we were asked to include a section addressing metabolic dysfunction-associated fatty liver disease (MAFLD), a new name and set of diagnosis criteria proposed by experts in the field in 2020\cite{Eslam2020-ud}. In 2023, another name and definition were introduced: metabolic dysfunction-associated steatotic liver disease (MASLD)\cite{Rinella2023-oe}. While the two new terms are still under debate\cite{Ramirez-Mejia2024-jo}, in just two years, the review I wrote became obsolete. I strongly support this trajectory toward precision medicine, as both MAFLD and MASLD are examples of reducing ambiguity, removing stigma from medical terms, and replacing invasive diagnostic methods with minimally invasive measurements, all of which benefit patients, researchers, and clinicians. Similarly, for metabolic syndrome, we are moving from using binary definitions\cite{Kraja2011-ja}, to severity indices\cite{Gurka2014-ib}, to multivariate genomic analysis\cite{Park2024-da}. While METSS served as a robust tool for this specific study, future applications will have better methods that reflect our deepening understanding of metabolic health.
    
    The reprogramming framework proposed in Chapter 3 offers significant potential for expanding our understanding of adipocytes and mechanisms underlying IR beyond the current models. However, translating this computational framework into a viable experimental system presents substantial technical hurdles. The multifactorial nature of cellular reprogramming requires rigorous optimization of experimental conditions to evaluate candidate transcription factors (TFs), adipocytes are particularly difficult to work with due to cellular fragility, and many TFs were proposed. More adipocyte and adipose tissue datasets would be beneficial. Ideally, the proposed framework could also be extended to model other understudied cell types with primarily transcriptomic data.
    
    During my PhD, I used functional genomics combined with human genetics to determine causal mechanisms of insulin resistance, with a focus on key regulators and cell types in adipose tissue.  Looking forward, the field should prioritize high-resolution, multimodal research on adipocyte biology to further elucidate the mechanisms driving insulin resistance. 

\end{dissertationepilogue}

\bibliographystyle{unsrtnat}
\bibliography{bib_4.bib}